\documentclass[10pt]{article}
\usepackage[T1]{fontenc}
\usepackage[utf8]{inputenc}
\usepackage[french]{babel}
\usepackage{amsmath,amssymb,mathrsfs,mathpazo}
\usepackage{geometry}
\geometry{a4paper,margin=2cm}
\usepackage{graphicx,array,tabularx,tabularray,colortbl,multirow,multicol,booktabs,enumitem,pifont,ulem,fancybox,diagbox,cancel}
\usepackage{pstricks,pstricks-add,pst-plot,pst-tree,pst-node,pst-text,pst-func,pst-math,pst-bspline,pst-3dplot,pst-eucl,pst-solides3d}
\let\pscalculate\undefined
\usepackage{pst-bezier}
\usepackage{tikz,pgfplots,tkz-tab}
\usetikzlibrary{arrows,arrows.meta,patterns,calc,positioning,trees,decorations.markings,shapes.geometric,shapes.misc,shapes.arrows,intersections,decorations.pathmorphing,decorations.pathreplacing,backgrounds,shadows,babel}
\pgfplotsset{compat=1.18}
\usepgfplotslibrary{fillbetween}
\usepackage[np]{numprint}
\usepackage{listings,makecell,pas-tableur,twoopt}
\let\barre\undefined
\usepackage[pstricks]{tablvar}
\definecolor{darkgreen}{RGB}{0,100,0}
\providecommand{\red}{\color{red}}
\providecommand{\blue}{\color{blue}}
\providecommand{\green}{\color{green}}
\providecommand{\black}{\color{black}}
\providecommand{\white}{\color{white}}
\providecommand{\gray}{\color{gray}}
\providecommand{\cyan}{\color{cyan}}
\providecommand{\magenta}{\color{magenta}}
\providecommand{\yellow}{\color{yellow}}
\providecommand{\violet}{\color{violet}}
\providecommand{\orange}{\color{orange}}
\providecommand{\euro}{\texteuro}
\providecommand{\EUR}{\texteuro}
\providecommand{\ieme}{\textsuperscript{e}}
\providecommand{\degres}{\ensuremath{^\circ}}
\providecommand{\degree}{\ensuremath{^\circ}}
\providecommand{\vv}[1]{\overrightarrow{#1}}
\providecommand{\Coord}[1]{#1}
\providecommand{\decosix}{}
\providecommand{\starredbullet}{\ensuremath{\bullet}}
\pagestyle{empty}
\setlength{\parindent}{0pt}
\usepackage{adjustbox}\DeclareUnicodeCharacter{2236}{:}\DeclareUnicodeCharacter{2190}{\ensuremath{\leftarrow}}
\providecommand{\R}{\mathbb{R}}
\providecommand{\N}{\mathbb{N}}
\providecommand{\Z}{\mathbb{Z}}
\providecommand{\Q}{\mathbb{Q}}
\providecommand{\D}{\mathbb{D}}
\providecommand{\C}{\mathbb{C}}
\providecommand{\e}{{\mathrm{e}}}
\providecommand{\vect}[1]{\overrightarrow{#1}}
\providecommand{\vectt}[1]{\overrightarrow{\mathrm{#1}}}
\providecommand{\vvt}[1]{\overrightarrow{\mathrm{#1}}}
\providecommand{\np}[2][]{{#2~#1}}
\providecommand{\vv}[1]{{\overrightarrow{#1}}}
\renewcommand{\d}{{\mathrm{d}}}
\providecommand{\degres}{{^\circ}}
\providecommand{\degre}{{^\circ}}
\providecommand{\blue}{\color{blue}}
\providecommand{\red}{\color{red}}
\providecommand{\enskip}{\quad}
\providecommand{\llbracket}{\lbrack\!\lbrack}
\providecommand{\rrbracket}{\rbrack\!\rbrack}
\providecommand{\euro}{{\text{€}}}
\providecommand{\no}{{\text{n°}}}
\providecommand{\interval}[2]{{[#1;#2]}}
\providecommand{\intervFO}[2]{{[#1;#2[}}
\providecommand{\dotfill}{\ldots}

\providecommand{\ds}{\displaystyle}
\providecommand{\textding}[1]{#1}
\providecommand{\Oij}{\(\mathrm{O};\vec{\imath},\vec{\jmath}\)}
\providecommand{\Oijk}{\(\mathrm{O};\vec{\imath},\vec{\jmath},\vec{k}\)}
\providecommand{\Ouv}{\(\mathrm{O};\vec{u},\vec{v}\)}
\providecommand{\intervFF}[2]{\ensuremath{[#1;#2]}}
\renewcommand{\euro}{\ifmmode\text{\texteuro}\else\texteuro\fi}
\DeclareUnicodeCharacter{00B0}{\ifmmode^{\circ}\else\textdegree\fi}
\DeclareUnicodeCharacter{20AC}{\ifmmode\text{\texteuro}\else\texteuro\fi}
% Énoncés : Éducation nationale. Transcription et corrigés : APMEP, Denis Vergès et contributeurs. Découpage et classement : CoopMaths, Sébastien Lozano, Éric Elter et contributeurs. Adaptation de la présentation : Maths Explorer.
% CC BY-SA : https://coopmaths.fr/www/about/
% https://coopmaths.fr/alea/static/bac/2023/tex/bac_2023_03_sujet2_ameriquenord_1.tex
\begin{document}
\noindent\textit{APMEP / CoopMaths. CC BY-SA. Présentation adaptée par Maths Explorer.}\par\medskip
\medskip

\textbf{\large Partie A}

\medskip

Le plan est muni d'un repère orthogonal.

On considère une fonction $f$ définie et dérivable sur $\mathbb{R}$. On note $f'$ sa fonction dérivée. On donne ci-dessous la courbe représentative de la fonction dérivée $f'$.

\begin{center}
\psset{unit=1cm,arrowsize=2pt 3}
\begin{pspicture*}(-7.1,-8)(3.1,3)
\psaxes[linewidth=1.25pt,labelFontSize=\scriptstyle]{->}(0,0)(-7,-7.95)(3,3)
\psgrid[gridlabels=0pt,subgriddiv=1,gridwidth=0.3pt]
\psplot[plotpoints=2000,linewidth=1.25pt,linecolor=red]{-7}{3}{x dup mul 3 x mul sub 1 add 2.71828 x exp mul}
\end{pspicture*}

\end{center}

Dans cette partie, les résultats seront obtenus par lecture graphique de la courbe représentative de la fonction dérivée $f'$. Aucune justification n'est demandée.

\medskip

\begin{enumerate}
\item Donner le sens de variation de la fonction $f$ sur $\mathbb{R}$. On utilisera des valeurs approchées si besoin.
\item Donner les intervalles sur lesquels la fonction $f$ semble être convexe.
\end{enumerate}

\medskip

\textbf{\large Partie B}

\medskip

On admet que la fonction $f$ de la partie $\mathbf{A}$ est définie sur $\mathbb{R}$ par 
\[f(x)=\left(x^{2}- 5 x + 6\right) \mathrm{e}^{x}.\]

On note $\mathcal{C}$ la courbe représentative de la fonction $f$ dans un repère.

\medskip

\begin{enumerate}
\item 
	\begin{enumerate}
		\item Déterminer la limite de la fonction $f$ en $+\infty$.
		\item Déterminer la limite de la fonction $f$ en $-\infty$.
	\end{enumerate}
\item Montrer que, pour tout réel $x$, on a $f'(x)=\left(x^{2}-3 x+1\right) \mathrm{e}^{x}$.
\item En déduire le sens de variation de la fonction $f$.
\item Déterminer l'équation réduite de la tangente $(\mathcal{T})$ à la courbe $\mathcal{C}$ au point d'abscisse 0.
\end{enumerate}

On admet que la fonction $f$ est deux fois dérivable sur $\mathbb{R}$. On note $f''$ la fonction dérivée seconde de la fonction $f$. On admet que, pour tout réel $x$, on a $f''(x) = (x + 1)(x - 2) \mathrm{e}^{x}$.

\begin{enumerate}[resume]
\item 
	\begin{enumerate}
		\item Étudier la convexité de la fonction $f$ sur $\mathbb{R}$.
		\item Montrer que, pour tout $x$ appartenant à l'intervalle $[-1~;~2]$, on a $f(x) \leqslant x + 6$.	
	\end{enumerate}
\end{enumerate}

\bigskip

\end{document}
